% Curated, human-readable FLT37 chapter.

\section{Background: irregular primes and Vandiver's Theorem III}

Kummer's criterion (proved in \cref{ch:kummer-final}) characterises regular
primes: $p$ is regular iff $p$ does not divide the numerator of any of the
Bernoulli numbers $B_2, B_4, \ldots, B_{p-3}$.  Kummer showed in 1850 that
$\mathrm{FLT}_p$ holds for every regular prime~$p$.  This already proves
$\mathrm{FLT}_p$ for the vast majority of primes; the smallest irregular
primes are $37, 59, 67, 101, 103, \ldots$.

For irregular primes a more delicate argument is needed.  Vandiver gave
several supplementary criteria; the relevant one here is:

\begin{quote}
  \textbf{Vandiver's Theorem III.}  Let $p$ be an irregular prime with
  $p \equiv 1 \pmod 4$.  If for every irregular Bernoulli index $2k$ of $p$
  (i.e.\ every $k$ in $[1, (p-3)/2]$ with $p \mid \mathrm{num}(B_{2k})$) the
  integer $k$ itself is even, then $\mathrm{FLT}_p$ holds.
\end{quote}

For $p = 37$ the unique irregular index is $2k = 32$, so $k = 16$, which is
even, and $37 \equiv 1 \pmod 4$.  So Vandiver's Theorem III applies at $37$,
and proving $\mathrm{FLT}_{37}$ reduces to formalising Vandiver III at
$\ell = 37$.

\begin{definition}
  \label{def:isIrregularIndex}
  \lean{KummerCriterion.FLT37.IsIrregularIndex}\leanok
  An \emph{irregular index of $\ell$} is a positive integer $k$ with
  $1 \le k$ and $2k \le \ell - 3$ such that
  $\ell \mid \mathrm{num}(B_{2k})$.
\end{definition}

\begin{definition}
  \label{def:vandiverIIIHypothesis}
  \lean{KummerCriterion.FLT37.VandiverIIIHypothesis}\leanok
  \uses{def:isIrregularIndex}
  The \emph{Vandiver III parity hypothesis} for $\ell$ asserts:
  $\ell \equiv 1 \pmod 4$, and every irregular index $k$ of $\ell$
  is even.
\end{definition}

\begin{theorem}
  \label{thm:vandiverIIIHypothesis_thirtyseven}
  \lean{KummerCriterion.FLT37.vandiverIIIHypothesis_thirtyseven}\leanok
  \uses{def:vandiverIIIHypothesis}
  The prime $\ell = 37$ satisfies the Vandiver III parity hypothesis.
\end{theorem}

\begin{proof}\leanok
  By direct computation: $37 \equiv 1 \pmod 4$, and the only irregular
  index of $37$ in the range $[1, 17]$ is $k = 16$, which is even.
\end{proof}

\section{The case-decomposition}

Following \texttt{flt-regular}'s organisation of Kummer's argument for
regular primes, we split Vandiver III into a first-case half and a
second-case half, defined after the \texttt{MayAssume.coprime} reduction
that pulls a common factor out of $(a,b,c)$.

\begin{definition}
  \label{def:vandiverIII}
  \lean{KummerCriterion.FLT37.VandiverIII}\leanok
  \uses{def:vandiverIIIHypothesis}
  $\mathrm{VandiverIII}$ asserts: every prime $\ell$ satisfying the parity
  hypothesis satisfies $\mathrm{FermatLastTheoremFor}\ \ell$.
\end{definition}

\begin{definition}
  \label{def:vandiverIIICaseI}
  \lean{KummerCriterion.FLT37.VandiverIIICaseI}\leanok
  \uses{def:vandiverIIIHypothesis}
  $\mathrm{VandiverIIICaseI}$ asserts: under the parity hypothesis, if
  $\ell \nmid abc$ then $a^\ell + b^\ell \neq c^\ell$.
\end{definition}

\begin{definition}
  \label{def:vandiverIIICaseII}
  \lean{KummerCriterion.FLT37.VandiverIIICaseII}\leanok
  \uses{def:vandiverIIIHypothesis}
  $\mathrm{VandiverIIICaseII}$ asserts: under the parity hypothesis, if
  $a,b,c$ are coprime and $\ell \mid abc$, then $a^\ell + b^\ell \neq c^\ell$.
\end{definition}

\begin{theorem}
  \label{thm:vandiverIII_of_caseI_caseII}
  \lean{KummerCriterion.FLT37.vandiverIII_of_caseI_caseII}\leanok
  \uses{def:vandiverIII, def:vandiverIIICaseI, def:vandiverIIICaseII}
  $\mathrm{VandiverIIICaseI} \land \mathrm{VandiverIIICaseII}
    \;\Longrightarrow\; \mathrm{VandiverIII}$.
\end{theorem}

\begin{proof}\leanok
  Reduce $\mathrm{FermatLastTheoremFor}\ \ell$ to the integer statement via
  $\mathrm{fermatLastTheoremFor\_iff\_int}$ and apply
  $\mathrm{MayAssume.coprime}$ to assume $\gcd(a,b,c) = 1$.  Then split on
  $\ell \mid abc$: case~II handles the divisibility branch, case~I the
  non-divisibility branch.
\end{proof}

\begin{theorem}
  \label{thm:fermatLastTheoremFor_thirtyseven_of_caseI_caseII}
  \lean{KummerCriterion.FLT37.fermatLastTheoremFor_thirtyseven_of_caseI_caseII}\leanok
  \uses{thm:vandiverIIIHypothesis_thirtyseven, thm:vandiverIII_of_caseI_caseII}
  $\mathrm{VandiverIIICaseI} \land \mathrm{VandiverIIICaseII}
    \;\Longrightarrow\; \mathrm{FermatLastTheoremFor}\ 37$.
\end{theorem}

\begin{proof}\leanok
  Specialise \cref{thm:vandiverIII_of_caseI_caseII} to $\ell = 37$ using the
  verified parity-hypothesis witness
  (\cref{thm:vandiverIIIHypothesis_thirtyseven}).
\end{proof}

This reduces $\mathrm{FLT}_{37}$ to two named hypotheses
$\mathrm{VandiverIIICaseI}$ and $\mathrm{VandiverIIICaseII}$ at $\ell = 37$.
The next two sections describe the substantive mathematical content of each.

\section{Case I: the Mirimanoff route}

For a putative case-I solution $a^p + b^p = c^p$ with $p \nmid abc$, set
$t = -a / b \pmod p$.  Classical Mirimanoff theory (Ribenboim, \emph{13
Lectures on Fermat's Last Theorem}, Lecture VIII) gives a chain of
congruences in $\mathbb{F}_p$ relating $t$ to the values of the
\emph{Mirimanoff polynomials}
\[
  \varphi_n(t) \;:=\; \sum_{j=1}^{p-1} j^{n-1}\, t^j \quad\in\quad
    \mathbb F_p[t],
  \qquad n = 1, 2, \ldots, p-2.
\]
These polynomials arise from the cyclotomic-unit factorisation of
$a^p + b^p$ combined with Stickelberger's annihilator and the
$p$-adic logarithm.

The Lean development captures the chain through three predicates of
increasing strength.

\begin{definition}
  \label{def:mirimanoffPolynomialVanishing}
  \lean{KummerCriterion.FLT37.MirimanoffPolynomialVanishing}\leanok
  $\mathrm{MirimanoffPolynomialVanishing}$ asserts: for every $n$ with
  $2 \le n \le p - 3$, $\varphi_n(t) \equiv 0 \pmod p$.
\end{definition}

\begin{definition}
  \label{def:mirimanoffBernoulliIdentity}
  \lean{KummerCriterion.FLT37.MirimanoffBernoulliIdentity}\leanok
  $\mathrm{MirimanoffBernoulliIdentity}$ asserts the product congruence
  $\varphi_n(t) \cdot B_{p-n} \equiv 0 \pmod p$ for every odd $n$ with
  $2 \le n \le p - 3$.
\end{definition}

\begin{definition}
  \label{def:mirimanoffBernoulliConclusion}
  \lean{KummerCriterion.FLT37.MirimanoffBernoulliConclusion}\leanok
  $\mathrm{MirimanoffBernoulliConclusion}$ asserts: for every $n$ with
  $2 \le n \le p - 3$ such that $t^n \not\equiv 1 \pmod p$,
  $p \mid \mathrm{num}(B_{p-n})$.
\end{definition}

The Mirimanoff--Bernoulli identity ($\mathrm{MBI}$) is the classical
Mirimanoff theorem.  Combined with the Vandiver III parity hypothesis it
forces $\varphi_n(t) \equiv 0 \pmod p$ for every $n \equiv 3 \pmod 4$ in
range, because the parity hypothesis rules out
$p \mid \mathrm{num}(B_{p-n})$ when $(p-n)/2$ is odd, so the other factor
of the product congruence must vanish.  That in turn yields the
Bernoulli-divisibility conclusion ($\mathrm{MBC}$).

\begin{theorem}
  \label{thm:caseI_thirtyseven_of_bernoulli_conclusion}
  \lean{KummerCriterion.FLT37.caseI_thirtyseven_of_bernoulli_conclusion}\leanok
  \uses{def:mirimanoffBernoulliConclusion, thm:vandiverIIIHypothesis_thirtyseven}
  At $\ell = 37$, $\mathrm{MirimanoffBernoulliConclusion}\ 37\ a\ b$ for a
  case-I solution $(a,b,c)$ produces a contradiction (i.e.\ no such
  solution exists).
\end{theorem}

This is the case-I closure: once $\mathrm{MBC}$ is discharged from
$\mathrm{MBI}$ (which is the genuine open analytic content of case~I),
case~I at $\ell = 37$ is proved.

\section{Case II: descent under $37 \nmid h^+$}

Case~II at irregular primes is substantially harder than case~I.  The
classical Washington 9.4 / Lehmer--Vandiver descent argument requires
the hypothesis $p \nmid h^+$ (Vandiver's conjecture for $p$), where
$h^+ = h(\mathbb Q(\zeta_p)^+)$ is the class number of the maximal real
subfield.

At $p = 37$ this hypothesis is \emph{not} an assumption: it is proved
unconditionally in the cyclotomic-units chapter of this blueprint, via
Sinnott's index formula combined with a direct Bernoulli-table
computation at the irregular index~$32$.

\begin{theorem}
  \label{thm:flt37_not_dvd_hPlus}
  \lean{KummerCriterion.FLT37.Sinnott.flt37_not_dvd_hPlus}\leanok
  $37 \nmid h^+(\mathbb Q(\zeta_{37}))$.
\end{theorem}

\begin{proof}\leanok
  Combines two ingredients.  First, Sinnott's index formula
  \[
    [\mathcal O_{\mathbb Q(\zeta_p)^+}^\times : C^+] \;=\; 2^{(p-3)/2} \cdot h^+
  \]
  (developed in \cref{ch:cyclotomic-units}) reduces $p \nmid h^+$ to the
  $p$-coprimality of the cyclotomic-unit index $[E^+ : C^+]$.  Second, the
  $p$-coprimality of that index follows from the non-divisibility
  $p \nmid B_{2k}$ at all relevant $k$ except possibly the irregular ones;
  at $p = 37$ the unique irregular index $k = 16$ is discharged by a
  direct congruence check on $B_{32}\bmod 37^2$ (the Bernoulli-table
  computation).
\end{proof}

Under $37 \nmid h^+$, the case-II descent is captured by the predicate
$\mathrm{VandiverLemma1Thirtyseven}$, which encodes the
regularity-free structural input needed for the Washington 9.4 descent.

\begin{definition}
  \label{def:vandiverLemma1Thirtyseven}
  \lean{KummerCriterion.FLT37.VandiverLemma1Thirtyseven}\leanok
  $\mathrm{VandiverLemma1Thirtyseven}$ asserts: for any coprime
  integers $a,b,c$ with $37 \mid abc$, $a^{37} + b^{37} \neq c^{37}$.
\end{definition}

\begin{theorem}
  \label{thm:caseII_thirtyseven_of_vandiverLemma1}
  \lean{KummerCriterion.FLT37.caseII_thirtyseven_of_vandiverLemma1}\leanok
  \uses{def:vandiverLemma1Thirtyseven, def:vandiverIIICaseII}
  $\mathrm{VandiverLemma1Thirtyseven}$ implies $\mathrm{VandiverIIICaseII}$
  at $\ell = 37$.
\end{theorem}

\subsection*{Where $\mathrm{VandiverLemma1Thirtyseven}$ comes from}

The Lehmer--Vandiver machinery in
\texttt{KummerCriterion/FLT37/LehmerVandiver/} formalises the
Washington 9.4 descent in the $\sigma$-stable product form recommended
by the 2026-05-27 expert review.  The descent proceeds as follows.

Starting from a putative case-II solution at level $m \ge 1$
(i.e.\ $x^p + y^p = \varepsilon \cdot ((\zeta-1)^{m+1} \cdot z)^p$ in
$\mathcal O_{\mathbb Q(\zeta_p)}$ with $\zeta - 1 \nmid y, z$), one
factors $x^p + y^p = \prod_\eta (x + y\eta)$ over the $p$-th roots of
unity.  The ideals $\mathfrak a(\eta)$ such that
$(\mathfrak a(\eta))^p$ is the $\mathfrak p$-coprime part of
$\bigl((x + y\eta)/(\zeta-1)\bigr)/(\mathfrak m)$ are pairwise coprime
$p$-th powers and provide the descent ideals.

Because the raw quotient $\mathfrak a(\eta)/\mathfrak a(\eta_0)$ is not
fixed by complex conjugation, the descent at the unit-class level is run
on the $\sigma$-stable product $\mathfrak a(\eta) \cdot
\mathfrak a(\eta^{-1})$ and its anchored ratio against
$\mathfrak a(\eta_0) \cdot \mathfrak a(\eta_0^{-1})$.  Under
$37 \nmid h^+$ this ratio is principal in $\mathcal O_{K^+}$, giving
\emph{real} descent generators $x, y \in \mathcal O_{K^+}$ with
$(x) \cdot J = (y) \cdot J_0$.  Combined with the $p$-th-power
identity coming from the polynomial factorisation, these generators
produce a Cramer-style Fermat-like equation
\[
  \varepsilon_1' \, X^{37} + \varepsilon_2' \, Y^{37} \;=\; Z^{37}
\]
in $\mathcal O_K$ with $X, Y, Z$ real (fixed by complex conjugation) and
$\varepsilon_1', \varepsilon_2' \in \mathcal O_K^\times$ explicitly
constructed units.  Iterating the descent yields strictly decreasing
$\mathfrak p$-adic exponents, hence a contradiction.

The substantive open content is closing this chain: the Cramer Fermat
identity is shipped at $\mathcal O_K$ level (file
\texttt{LehmerVandiver/CaseII/ProductDescent.lean}); discharging
$\mathrm{VandiverLemma1Thirtyseven}$ then requires using this identity
to produce a new $\mathrm{RealCaseIIData37}$ at a strictly smaller
$\mathfrak p$-adic exponent, completing the infinite descent.  This is
the only remaining mathematical input.

\section{Assembly}

\begin{theorem}
  \label{thm:fermatLastTheoremFor_thirtyseven_of_conclusion_and_vandiverLemma1}
  \lean{KummerCriterion.FLT37.fermatLastTheoremFor_thirtyseven_of_conclusion_and_vandiverLemma1}\leanok
  \uses{thm:caseI_thirtyseven_of_bernoulli_conclusion,
        thm:caseII_thirtyseven_of_vandiverLemma1,
        thm:fermatLastTheoremFor_thirtyseven_of_caseI_caseII,
        thm:flt37_not_dvd_hPlus}
  If
  \begin{enumerate}
    \item every case-I solution $(a,b,c)$ at $\ell = 37$ satisfies
      $\mathrm{MirimanoffBernoulliConclusion}\ 37\ a\ b$, and
    \item $\mathrm{VandiverLemma1Thirtyseven}$ holds,
  \end{enumerate}
  then $\mathrm{FermatLastTheoremFor}\ 37$.
\end{theorem}

\section{Supporting machinery developed in the project}

This section records the principal named theorems developed in the
project — \emph{including} pieces that were not used on the critical path
to the FLT37 statement above.  They are part of the project's
deliverable formalisation of cyclotomic-field arithmetic and are
referenced from multiple chapters of the wider blueprint.

\subsection*{Sinnott's index formula}

\begin{theorem}
  \label{thm:sinnottIndexFormula_of_regulatorIdentity}
  \lean{KummerCriterion.FLT37.Sinnott.sinnottIndexFormula_of_regulatorIdentity}\leanok
  Under the Kummer--Dirichlet regulator identity (the analytic input,
  proved via the deleted-Fourier formula), Sinnott's index formula holds:
  \[
    [\mathcal O_{\mathbb Q(\zeta_p)^+}^\times : C^+]
      \;=\; 2^{(p-3)/2}\, h^+.
  \]
\end{theorem}

\begin{theorem}
  \label{thm:cyclotomicUnitIndex_primeConductor_pPrimary_of_sinnottIndexFormula}
  \lean{KummerCriterion.cyclotomicUnitIndex_primeConductor_pPrimary_of_sinnottIndexFormula}\leanok
  \uses{thm:sinnottIndexFormula_of_regulatorIdentity}
  Sinnott's formula transports the divisibility $p \mid h^+$ to the
  cyclotomic-unit index $p \mid [E^+ : C^+]$ (and conversely), giving
  the unit-side characterisation used to bypass Vandiver's conjecture
  at $p = 37$ via a direct Bernoulli-table check.
\end{theorem}

\subsection*{Thaine's theorem and circular units}

\begin{theorem}
  \label{thm:circularSubgroupKplus_eq_sinnott_eq_washington}
  \lean{KummerCriterion.Thaine.circularSubgroupKplus_eq_sinnott_eq_washington}\leanok
  Sinnott's and Washington's two conventions for the group of
  \emph{circular units} of $\mathbb Q(\zeta_p)^+$ coincide.  The
  identification underwrites the use of either convention in
  downstream Thaine-style annihilation arguments.
\end{theorem}

\begin{theorem}
  \label{thm:thaineAuxiliaryExistence_of_prime}
  \lean{KummerCriterion.Thaine.thaineAuxiliaryExistence_of_prime}\leanok
  For every prime $\ell$ in the Chebotarev-positive density set
  $\{\ell : \ell \equiv 1 \pmod p, \ell \text{ totally split in some
  auxiliary field}\}$, the \emph{Thaine auxiliary} construction
  produces an annihilator of the class group $\mathrm{Cl}(\mathbb Q(\zeta_p)^+)$
  drawn from a circular unit at $\ell$.  Combined with a finite-set
  exclusion, this gives existence of Thaine auxiliaries outside any
  fixed finite set of bad primes.
\end{theorem}

\begin{theorem}
  \label{thm:infinite_setOf_thaineAuxiliary}
  \lean{KummerCriterion.Thaine.infinite_setOf_thaineAuxiliary}\leanok
  \uses{thm:thaineAuxiliaryExistence_of_prime}
  The set of Thaine auxiliary primes is infinite.
\end{theorem}

\subsection*{Reflection / Spiegelungssatz}

\begin{theorem}
  \label{thm:weakReflection_dvd_hMinus_of_dvd_hPlus_units}
  \lean{KummerCriterion.weakReflection_dvd_hMinus_of_dvd_hPlus_units}\leanok
  The \emph{weak reflection principle} on the unit side: divisibility
  of the plus class number $p \mid h^+$ implies a non-trivial
  $\epsilon_i$-eigencomponent in the minus class group for some
  reflection-component index $i$.  This is the project's bridge from
  unit-side data to class-side data, used as one input to the
  Lehmer--Vandiver descent.
\end{theorem}

\begin{theorem}
  \label{thm:dvd_h_iff_exists_dvd_bernoulli_of_weakReflection}
  \lean{KummerCriterion.dvd_h_iff_exists_dvd_bernoulli_of_weakReflection}\leanok
  \uses{thm:weakReflection_dvd_hMinus_of_dvd_hPlus_units}
  Conditional on the weak reflection principle, $p \mid h$ is
  equivalent to $p$ dividing a Bernoulli numerator in the
  Kummer range — an alternative form of Kummer's criterion routed
  through the reflection / Spiegelungssatz machinery rather than the
  Sinnott / circular-unit route.
\end{theorem}

\subsection*{The CM splitting $h^+ \mid h$}

\begin{theorem}
  \label{thm:hPlus_dvd_h}
  \lean{KummerCriterion.hPlus_dvd_h}\leanok
  In the CM field $\mathbb Q(\zeta_p)/\mathbb Q(\zeta_p)^+$, the plus
  class number divides the full class number,
  $h(\mathbb Q(\zeta_p)^+) \mid h(\mathbb Q(\zeta_p))$.  This is
  proved unconditionally via faithful flatness of
  $\mathcal O_{\mathbb Q(\zeta_p)}/\mathcal O_{\mathbb Q(\zeta_p)^+}$
  and the injectivity of the induced map of class groups.
\end{theorem}

\section{Outlook}

Modulo the unconditional $37 \nmid h^+$
(\cref{thm:flt37_not_dvd_hPlus}), the proof of
$\mathrm{FLT}_{37}$ in this project reduces to two named hypotheses:

\begin{enumerate}
  \item \textbf{Case I:} the per-solution Bernoulli-divisibility
  predicate $\mathrm{MirimanoffBernoulliConclusion}\ 37$, derived
  from Mirimanoff's classical theorem ($\mathrm{MBI}$) combined with the
  verified parity hypothesis.
  \item \textbf{Case II:} $\mathrm{VandiverLemma1Thirtyseven}$, encoding
  the Washington 9.4 / Lehmer--Vandiver descent under $37 \nmid h^+$ on
  $\sigma$-stable real data.
\end{enumerate}

Both are explicit cyclotomic-integer statements at the single prime
$37$.  Neither requires a fresh layer of class field theory beyond what
\texttt{flt-regular} and the cyclotomic-units chapter already supply.
